% Converted from README.md
\subsection{Java - JUnit, JaCoCo}\label{java---junit-jacoco}

\href{https://junit.org/}{JUnit} je jedan od najpopularnijih radnih okvira za testiranje jedinica koda u programskom jeziku Java. Neke od osobina JUnit radnog okvira su jednostavnost pisanja testova uz bogat skup anotacija koje opisuju testove, kao i veoma velika podrška za najčešće situacije u procesu pisanja testova kao što su uslovno uključivanje/isključivanje testova na osnovu promenljivih iz okruženja, operativnog sistema, proizvoljnih predikata itd.

JUnit svoje artefakte isporučuje na \href{https://search.maven.org/}{Maven Central} i moguće ga je uključiti u Maven (i Maven-kompatibilne) alate za prevođenje. U ovom primeru ćemo koristiti drugi popularni alat za prevođenje pod imenom \href{https://gradle.org/}{Gradle}. Gradle koristi datoteku \texttt{build.gradle} (slično kao što Maven koristi \texttt{pom.xml}) za definiciju projekta, zavisnosti i podešavanja dodataka.

Kreirajmo novi projekat:

\begin{minted}{bash}
$ gradle init
\end{minted}

Možemo u \texttt{build.gradle} dodati zavisnost za JUnit:

\begin{minted}{text}
  | dependencies {
  |    ...
+ |    testImplementation "junit:junit:4.13"
  | }
\end{minted}

Za praćenje pokrivenosti koda, koristićemo \href{https://www.jacoco.org/jacoco/}{JaCoCo}. JaCoCo možemo lako uključiti u projekat dodavanjem niske \texttt{\textquotesingle{}jacoco\textquotesingle{}} u spisak dodataka za projekat i dodatno konfigurisati izveštaj koji JaCoCo pravi:

\begin{minted}{text}
  | plugins {
  |     id 'java'
+ |     id 'jacoco'
  | }
  |
  | ...
  |
+ | jacocoTestReport {
+ |    reports {
+ |        xml.required = false
+ |        csv.required = false
+ |        html.outputLocation = layout.buildDirectory.dir('jacocoHtml')
+ |    }
+ | }
+ |
+ | check.dependsOn jacocoTestReport
\end{minted}

Nakon ovoga možemo iskoristiti sledeće komande:

\begin{itemize}
\tightlist
\item
  \texttt{gradle\ build} - prevođenje projekta
\item
  \texttt{gradle\ test} - pokretanje testova jedinica koda
\item
  \texttt{gradle\ jacocoTestReport} - kreiranje izveštaja o pokrivenosti
\end{itemize}

Izveštaji o pokrenutim testovima i statusu izvršavanja mogu se naći u direktorijumu \texttt{build/reports/tests/test} u HTML formatu. JaCoCo izveštaj o pokrivenosti može se naći u direktorijumu koji smo naveli u datoteci \texttt{build.gradle} - \texttt{build/jacocoHtml}, takođe u HTML formatu.

Pregled nekih od korisnih anotacija JUnit radnog okvira:

\begin{itemize}
\tightlist
\item
  \texttt{@Test} - Radni okvir će pokrenuti ovaj metod automatski prilikom pokretanja testova.
\item
  \texttt{@TestFactory} - Metod koji generiše testove u vremenu izvršavanja. Najčešće se koristi da pokrene nasumične testove ili testove bazirane na spoljnim podacima.
\item
  \texttt{@DisplayName} - Čini izveštaje čitljivijim tako što testovima daje navedeno ime.
\item
  \texttt{@BeforeAll}/\texttt{@BeforeEach} - Izvršava metod pre svih odnosno svakog testa.
\item
  \texttt{@AfterAll}/\texttt{@AfterEach} - Izvršava metod posle svih odnosno svakog testa.
\item
  \texttt{@Tag} - Dodaje oznaku testu radi kategorisanja testova u svite, npr. \texttt{@Tag("fast")} dodaje test u svitu sa oznakom \texttt{"fast"}.
\item
  \texttt{@Disabled} - Isključuje test metod iz radnog okvira.
\item
  \texttt{@Nested} - Koristi se u unutrašnjim klasama najčešće radi definisanja redosleda kojim se pokreću testovi.
\end{itemize}

Primeri korišćenja JUnit anotacija, test metoda i metoda pretpostavki:

\begin{minted}{java}
import org.junit.jupiter.api.*;
public class AppTest {
    @BeforeAll
    static void setup(){
        System.out.println("Executes a method Before all tests");
    }
    @BeforeEach
    void setupThis(){
        System.out.println("Executed Before each @Test method " +
                           "in the current test class");
    }
    @AfterEach
    void tearThis(){
        System.out.println("Executed After each @Test method " +
                           "in the current test class");
    }
    @AfterAll
    static void tear(){
        System.out.println("Executes a method After all tests");
    }
}
\end{minted}

\begin{minted}{java}
Assertions.assertAll("heading",
    () -> assertTrue(true),
    () -> assertEquals("expected", objectUnderTest.getSomething()
);
\end{minted}

\begin{minted}{java}
@TestFactory
Stream dynamicTests(MyContext ctx) {
    // Generates tests for every line in the file
    return Files.lines(ctx.testDataFilePath)
                .map(l -> dynamicTest("Test: " + l,
                                      () -> assertTrue(runTest(l))
                );
}
\end{minted}

\begin{minted}{java}
@Test
void exampleTest() {
    Assertions.assertTrue(trueBool);
    Assertions.assertFalse(falseBool);
    Assertions.assertNotNull(notNullString);
    Assertions.assertNull(notNullString);
    Assertions.assertNotSame(originalObject, otherObject);
    Assertions.assertEquals(4, 4);
    Assertions.assertNotEquals(3, 2);
    Assertions.assertArrayEquals(
        new int[] { 1, 2, 3 },
        new int[] { 1, 2, 3 },
        "Array Equal Test"
    );
    Iterable<Integer> listOne = new ArrayList<>(Arrays.asList(1,2,3,4));
    Iterable<Integer> listTwo = new ArrayList<>(Arrays.asList(1,2,3,4));
    Assertions.assertIterableEquals(listOne, listTwo);
    Assertions.assertTimeout(Duration.ofMillis(100), () -> {
        Thread.sleep(50);
        return "result";
    });
    Throwable exception = Assertions.assertThrows(
        IllegalArgumentException.class,
        () -> throw new IllegalArgumentException("error message");
    );
    Assertions.fail("not found good reason to pass");
}
\end{minted}

\begin{minted}{java}
@Test
void testAssumption() {
    System.setProperty("prop", "foo");
    Assumptions.assumeTrue("foo".equals(System.getProperty("prop")));
}
\end{minted}

\begin{minted}{java}
@Test
@EnabledForJreRange(min = JRE.JAVA_8, max = JRE.JAVA_11)
public void test1()
{
    System.out.println("Will run only on JRE between 8 and 11");
}

@Test
@EnabledOnJre({JRE.JAVA_8, JRE.JAVA_11})
public void test2()
{
    System.out.println("Will run only on JRE 8 and 11");
}

@Test
@DisabledForJreRange(min = JRE.JAVA_8, max = JRE.JAVA_11)
public void test3()
{
    System.out.println("Will NOT run on JRE between 8 and 11");
}

@Test
@DisabledOnJre({JRE.JAVA_8, JRE.JAVA_11})
public void test4()
{
    System.out.println("Will NOT run on JRE 8 and 11");
}
\end{minted}

\begin{minted}{java}
@Test
@EnabledOnOs({OS.LINUX, OS.WINDOWS})
void onLinuxOrWindows() {
    System.out.println("Will run on Linux or Windows.");
}

@Test
@DisabledOnOs({OS.WINDOWS, OS.SOLARIS, OS.MAC})
void notOnWindowsOrSolarisOrMac() {
    System.out.println("Won't run on Windows, Solaris or MAC!");
}
\end{minted}

\begin{minted}{java}
@Test
@EnabledIf("myCustomPredicate")
void enabled() {
    assertTrue(true);
}

@Test
@DisabledIf("myCustomPredicate")
void disabled() {
    assertTrue(true);
}

boolean myCustomPredicate() {
    return true;
}
\end{minted}

\begin{minted}{java}
@Test
@EnabledIfEnvironmentVariable(named = "ENV", matches = ".*oracle.*")
public void executeOnlyInDevEnvironment() {
    return true;
}

@Test
@DisabledIfEnvironmentVariable(named = "ENV", matches = ".*mysql.*")
public void disabledOnProdEnvironment() {
    return true;
}

@Test
@EnabledIfSystemProperty(named = "my.property", matches = "prod*")
public void onlyIfMyPropertyStartsWithProd() {
    return true;
}
\end{minted}
